% ===========================================================================
% CHAPTER 8
% ===========================================================================

\chapter{Evaluation, Security, and Limitations}
\label{chap:verification-security-discussion}

\section{Evaluation Method and Verification Results}
\label{sec:verification-method}

Evaluation combined structural inspection, static analysis, interface
traceability, and review of the available experimental outputs. Structural
inspection examined whether acquisition, dataset maintenance, analytics,
application services, and presentation components used consistent data
contracts. Static checks assessed Python syntax, TypeScript type consistency,
and frontend lint status. These methods establish structural confidence but do
not replace full integration, load, recovery, or penetration testing.

Table~\ref{tab:verification-results} summarizes the principal verification
results.

\begin{table}[H]
  \centering
  \small
  \setlength{\tabcolsep}{4pt}
  \caption{Verification results}
  \label{tab:verification-results}
  \begin{tabularx}{\textwidth}{@{}>{\RaggedRight\arraybackslash}p{0.25\textwidth}>{\RaggedRight\arraybackslash}p{0.29\textwidth}Y@{}}
    \toprule
    \textbf{Check} & \textbf{Result} & \textbf{Interpretation} \\
    \midrule
    Python syntax analysis & All 105 first-party Python files parsed successfully & Establishes syntax validity, not runtime behavior \\
    TypeScript compilation & Completed successfully without output generation & Establishes type consistency for the frontend project \\
    Frontend lint analysis & Reported 22 errors and 45 warnings & Identifies remaining code-quality issues despite successful type checking \\
    Backend runtime check & Not completed in the available environment & Full backend execution and endpoint behavior remain to be validated \\
    Automated test inventory & No substantive end-to-end regression suite was available & Functional, recovery, and concurrency claims require additional testing \\
    Interface and data tracing & Major acquisition, analytics, and dashboard paths were structurally consistent & Confirms component relationships but not every runtime outcome \\
    \bottomrule
  \end{tabularx}
\end{table}

The verification identified several secondary contracts that require further
alignment, including an auxiliary export path, a duplicate scope-switch path,
and a shared-key timeline query. These paths do not alter the principal
acquisition, analytical, or dashboard workflows, but they should be corrected
and covered by integration tests before broader deployment.

\section{Experimental IP-Crawler Results}
\label{sec:ip-experiment-results}

The experimental IP crawler produced separate summaries for runs without
Server Name Indication and with Server Name Indication. Each recorded run
contains 35,558 outcomes. Table~\ref{tab:ip-experiment-record} presents the
available results.

\begin{table}[H]
  \centering
  \small
  \setlength{\tabcolsep}{6pt}
  \caption{Experimental IP-crawler outcome summary}
  \label{tab:ip-experiment-record}
  \begin{tabularx}{\textwidth}{@{}Y>{\Centering\arraybackslash}p{0.20\textwidth}>{\Centering\arraybackslash}p{0.20\textwidth}@{}}
    \toprule
    \textbf{Outcome category} & \textbf{Without SNI} & \textbf{With SNI} \\
    \midrule
    SSL certificates found & 383 & 554 \\
    No service on port 443 & 404 & 558 \\
    TCP timeouts & 34,347 & 33,724 \\
    TLS timeouts & 5 & 26 \\
    SSL handshake errors & 22 & 31 \\
    Certificate parser errors & 0 & 0 \\
    Other errors & 397 & 665 \\
    \midrule
    Total recorded outcomes & 35,558 & 35,558 \\
    \bottomrule
  \end{tabularx}
\end{table}

The SNI-labelled run recorded 554 successful certificates, compared with 383
without SNI, a descriptive difference of 171. The available summary does not
include repeated trials, a complete environment description, per-address
results, or evidence that every other condition was held constant. The values
therefore demonstrate recorded crawler behavior but do not establish a causal
effect of SNI or a population-level conclusion about Pakistan or global TLS
deployment.

\section{Security and Scalability Assessment}
\label{sec:security-analysis}

Security analysis distinguishes certificate measurement from controls expected
in a deployed multi-user service. During collection, trust and hostname
validation are disabled so that presented certificates can be observed; this
activity must remain isolated from normal trusted-client behavior. The
dashboard is designed for controlled operation and would require
authentication, authorization, protected administrative actions, externalized
secrets, restricted service exposure, and hardened environment-specific
settings before broader deployment.

Scope is currently maintained as shared service state, so concurrent multi-user
operation requires request-local scope propagation. Internet-facing collection
also requires an approved operating policy covering authorization, rate
limits, opt-out handling, data retention, and responsible disclosure.

The system includes several mechanisms intended to support large datasets:
concurrent acquisition, atomic task claims, batch writes, persisted
checkpoints, compound indexes, materialized analytics, bounded candidate
queries, pagination, optional response caching, and client-side request reuse.
These mechanisms provide a scalability-oriented design, but their quantitative
effect has not been established through a controlled throughput, latency,
resource, or multi-user benchmark.

\section{System Limitations and Threats to Validity}
\label{sec:implementation-limitations-drift}

The main limitations of the developed platform are:

\begin{enumerate}[label=\arabic*.,leftmargin=0.42in]
  \item \textbf{Configuration consistency.} All components must use one validated dataset and service configuration to avoid inconsistent analytical results.
  \item \textbf{Dataset-update correctness.} Renewal checkpoints and temporary staging data require coordinated lifecycle management, while fingerprint deduplication assumptions should be enforced consistently.
  \item \textbf{Analytical calibration.} CA rankings, signature-strength values, shared-key exposure labels, and risk scores are project-defined and require external validation before being interpreted as scientific security measurements.
  \item \textbf{Geographic interpretation.} Country views are based on domain suffix and configured scope rather than physical server, subject, or issuer location.
  \item \textbf{Evaluation coverage.} Full backend integration, concurrency, recovery, usability, security, and performance evaluation remains incomplete.
  \item \textbf{Deployment readiness.} Production configuration, automated deployment, dependency locking, service supervision, and comprehensive test automation require further work.
\end{enumerate}

\noindent\textbf{Internal validity.}
Static and structural checks cannot establish successful behavior under every
runtime, concurrency, or interruption condition.

\noindent\textbf{Construct validity.}
The analytical scores operationalize project-defined rules. Shared public keys
do not prove compromise, encoded expiry does not establish stale authorization,
and scheduled expiration counts are not predictions.

\noindent\textbf{External validity.}
The domain and IP inputs are bounded datasets rather than random or complete
samples of global TLS deployment. The experimental IP results are separate from
the primary dashboard dataset.

\noindent\textbf{Reproducibility and temporal validity.}
Repeatable evaluation requires frozen datasets, complete environment
definitions, versioned native tools, recorded commands, and dated outputs.
Certificates, CT observations, network conditions, and domain ownership also
change over time, so analytical results represent the state captured by a
particular acquisition and refresh cycle.
